% ============================================================
% APEP 论文 LaTeX 模板 — paper_template.tex
% 基于 AER/AEJ 风格的经济学论文模板
% ============================================================

\documentclass[12pt]{article}

% --- 页面设置 ---
\usepackage[margin=1in]{geometry}
\usepackage{setspace}
\doublespacing

% --- 数学与符号 ---
\usepackage{amsmath, amssymb, amsfonts}

% --- 表格 ---
\usepackage{booktabs}
\usepackage{tabularx}
\usepackage{threeparttable}

% --- 图片 ---
\usepackage{graphicx}
\usepackage{float}
\usepackage{subcaption}

% --- 引用 ---
\usepackage{natbib}
\bibliographystyle{aer}  % AER 引用风格

% --- 链接 ---
\usepackage{hyperref}
\hypersetup{
    colorlinks=true,
    linkcolor=blue,
    citecolor=blue,
    urlcolor=blue
}

% --- 其他 ---
\usepackage{appendix}
\usepackage{enumitem}

% ============================================================
% 论文信息 — 请替换以下内容
% ============================================================
\title{%
  {PAPER_TITLE} \\
  \large Evidence from {POLICY_CONTEXT}%
}
\author{%
  APEP Working Paper No. apep\_XXXX \\
  Autonomous Policy Evaluation Project%
}
\date{\today}

% ============================================================
\begin{document}

\maketitle

% --- 摘要 ---
\begin{abstract}
\noindent
{ABSTRACT_TEXT: 描述研究问题、数据来源、识别策略、主要发现和政策含义。250 词以内。}

\bigskip
\noindent\textbf{Keywords:} {keyword1}, {keyword2}, {keyword3}

\noindent\textbf{JEL Codes:} {J01}, {H75}, {I38}
\end{abstract}

\newpage

% ============================================================
\section{Introduction}
% ============================================================

{INTRO_PARAGRAPH_1: 引出研究问题，说明政策背景的重要性。}

{INTRO_PARAGRAPH_2: 文献空白与本文贡献。}

This paper makes three contributions to the literature. First, ...
Second, ...
Third, ...

{INTRO_PARAGRAPH_3: 方法概述和主要发现预览。}

The remainder of this paper is organized as follows. Section~\ref{sec:background} provides institutional background. Section~\ref{sec:data} describes the data. Section~\ref{sec:strategy} presents the empirical strategy. Section~\ref{sec:results} reports the main results. Section~\ref{sec:robustness} discusses robustness checks. Section~\ref{sec:conclusion} concludes.

% ============================================================
\section{Institutional Background} \label{sec:background}
% ============================================================

{BACKGROUND: 政策的详细时间线、实施内容、覆盖范围。}

% 政策实施时间线表格
\begin{table}[H]
\centering
\caption{Policy Adoption Timeline}
\label{tab:timeline}
\begin{threeparttable}
\begin{tabular}{lcc}
\toprule
State & Adoption Year & Key Provisions \\
\midrule
{State1} & {Year1} & {Provision1} \\
{State2} & {Year2} & {Provision2} \\
\bottomrule
\end{tabular}
\begin{tablenotes}
\small
\item Notes: {表格注释}
\end{tablenotes}
\end{threeparttable}
\end{table}

% ============================================================
\section{Data} \label{sec:data}
% ============================================================

\subsection{Data Sources}
{DATA_SOURCES: 描述数据来源、API、时间覆盖、地理覆盖。}

\subsection{Variable Definitions}
{VARIABLES: 关键变量的定义和构造方法。}

\subsection{Summary Statistics}

\begin{table}[H]
\centering
\caption{Summary Statistics}
\label{tab:summary}
\begin{threeparttable}
\begin{tabular}{lccccc}
\toprule
Variable & N & Mean & SD & Min & Max \\
\midrule
{var1} & & & & & \\
{var2} & & & & & \\
\bottomrule
\end{tabular}
\begin{tablenotes}
\small
\item Notes: {来源和样本描述}
\end{tablenotes}
\end{threeparttable}
\end{table}

% ============================================================
\section{Empirical Strategy} \label{sec:strategy}
% ============================================================

% === DiD 模型规范 ===
We exploit the staggered adoption of {POLICY_NAME} across states using a difference-in-differences framework. Our baseline specification is:

\begin{equation}
Y_{st} = \alpha + \beta \cdot D_{st} + \gamma_s + \delta_t + \epsilon_{st}
\label{eq:main}
\end{equation}

where $Y_{st}$ is {OUTCOME_DESCRIPTION} in state $s$ in year $t$, $D_{st}$ is an indicator equal to one if state $s$ has adopted {POLICY_NAME} by year $t$, $\gamma_s$ and $\delta_t$ are state and year fixed effects, respectively.

% === RDD 模型规范（替代） ===
% We employ a sharp regression discontinuity design:
% \begin{equation}
% Y_i = \alpha + \tau \cdot \mathbb{1}[X_i \geq c] + f(X_i - c) + \epsilon_i
% \end{equation}

\subsection{Identification Assumptions}
{ASSUMPTIONS: 平行趋势假设的论证。}

\subsection{Modern DiD Estimators}
To address concerns about heterogeneous treatment effects in staggered designs \citep{goodman-bacon2021}, we supplement our TWFE estimates with:
\begin{itemize}
    \item The interaction-weighted estimator of \citet{sun-abraham2021}
    \item The group-time ATT estimator of \citet{callaway-santanna2021}
    \item The two-stage estimator of \citet{gardner2022}
\end{itemize}

% ============================================================
\section{Results} \label{sec:results}
% ============================================================

\subsection{Main Results}
\input{tables/main_results.tex}

{RESULTS_DISCUSSION: 解读主要回归结果，强调经济意义。}

\subsection{Event Study}
\begin{figure}[H]
\centering
\includegraphics[width=0.9\textwidth]{figures/event_study.pdf}
\caption{Event Study Estimates}
\label{fig:event_study}
\floatfoot{Notes: {图表注释}}
\end{figure}

{EVENT_STUDY_DISCUSSION: 事件研究图的解读，预趋势分析。}

% ============================================================
\section{Robustness} \label{sec:robustness}
% ============================================================

\subsection{Placebo Tests}
{PLACEBO: 安慰剂检验结果。}

\subsection{Alternative Specifications}
{ALT_SPECS: 不同设定的结果。}

\subsection{Heterogeneity Analysis}
{HETEROGENEITY: 子样本异质性分析。}

% ============================================================
\section{Conclusion} \label{sec:conclusion}
% ============================================================

{CONCLUSION: 总结主要发现、政策含义、局限性和未来研究方向。}

% ============================================================
% 参考文献
% ============================================================
\newpage
\bibliography{references}

% ============================================================
% 附录
% ============================================================
\newpage
\appendix
\section{Additional Results}
{APPENDIX_TABLES_FIGURES}

\end{document}
